\section{Related Work}\label{sec:related}
\subsection{The sim-to-real gap and domain randomisation}
The gap between simulated and real robot behaviour is a long-standing concern. Early work
in evolutionary robotics framed it as the reality gap and argued that controllers
transfer only when the simulator captures the right sources of variation, with noise
playing a central role in preventing overfitting to simulator artefacts
\cite{jakobi1995noise}. Domain randomisation generalises this idea by sampling simulator
parameters such as textures, lighting, and dynamics, so that the real world appears as
one more draw from a broad training distribution \cite{tobin2017domain}. Such methods are
powerful for learned systems but are typically applied as a blanket strategy; they do not
attribute the residual gap to a specific physical cause. This thesis takes the
complementary, analytic route of isolating one cause, ToF depth noise, and measuring its
effect directly.

\subsection{ToF and depth-sensor noise modelling}
Depth sensors have well-studied, sensor-specific error characteristics. Structured-light
and ToF cameras differ markedly, and ToF noise in particular is shaped by integration
time, multipath, and surface reflectance \cite{sarbolandi2015kinect,foix2011lockin}. For
consumer depth cameras, a widely used axial-and-lateral model expresses the depth
standard deviation as a function of range and viewing angle \cite{nguyen2012modeling},
while empirical characterisations of Kinect-class sensors report range-dependent
precision and systematic bias \cite{khoshelham2012accuracy}. A recurring finding is that
a large fraction of usable error concentrates at depth discontinuities as flying or mixed
pixels \cite{foix2011lockin,sarbolandi2015kinect}. We adopt these models as the expected
behaviour that our own measurements test, and we report where they do and do not hold for
the OAK-D SR PoE.

\subsection{Tabletop perception and pose estimation}
Classical tabletop perception pipelines segment objects from a dominant support plane,
often by fitting the plane with RANSAC and clustering the remaining points
\cite{rusu2008towards}. When the support surface is known and fronto-parallel, as in our
fixed overhead setup, the plane fit can be replaced by a simpler depth threshold; we make
this substitution deliberately and cite the RANSAC-based pipeline only as a point of
contrast. For evaluating pose, the BOP benchmark consolidated metrics and datasets for 6D
object pose \cite{hodan2018bop}; we use only its notion of aligning the visible top
surface, because our estimator produces a translation-only (3-DoF) pose rather than a
full 6D pose with orientation.

\subsection{Fiducial markers for ground truth}
Establishing a real-world reference frame requires reliable fiducials. Square markers such
as ArUco provide robust, occlusion-tolerant detection \cite{garridojurado2014aruco}, and
fast detectors improve corner localisation and detection speed
\cite{romeroramirez2018charuco}. We use a printed ChArUco board (ArUco markers combined
with a chessboard) to recover the camera-to-table transform by direct \texttt{solvePnP},
and individual markers to obtain per-object ground-truth positions.

\subsection{Grasping}
Data-driven grasp synthesis spans analytic and learned approaches and is surveyed
comprehensively in \citeA{bohg2014grasp}; learned multi-view grasping in clutter is an
active line within the host lab \cite{kasaei2021mvgrasp}. Our work does not contribute a
grasping method. Consistent with the proposal, we use a fixed, predefined grasp for known
objects so that any grasp error is attributable to perception rather than to grasp
planning, and we use this baseline purely to test whether pose-level errors propagate to
pick success.

% =====================================================================
